\thesisAppendix{А}{Скриншоты iOS-приложения}

В настоящем приложении приведены скриншоты основных пользовательских экранов iOS-приложения <<ВКурсе>>. Изображения собраны в композитные рисунки по тематическим разделам: каждый рисунок состоит из четырёх панелей --- (а), (б), (в) и (г), --- размещённых сеткой $2 \times 2$.

\begin{figure}[h]
  \centering
  \begin{minipage}[t]{0.45\textwidth}
    \centering
    \includegraphics[width=0.6\linewidth]{app-ios-onboarding-a.png}\\[3pt]
    \small (а)
  \end{minipage}\hfill
  \begin{minipage}[t]{0.45\textwidth}
    \centering
    \includegraphics[width=0.6\linewidth]{app-ios-onboarding-b.png}\\[3pt]
    \small (б)
  \end{minipage}\\[10pt]
  \begin{minipage}[t]{0.45\textwidth}
    \centering
    \includegraphics[width=0.6\linewidth]{app-ios-onboarding-c.png}\\[3pt]
    \small (в)
  \end{minipage}\hfill
  \begin{minipage}[t]{0.45\textwidth}
    \centering
    \includegraphics[width=0.6\linewidth]{app-ios-onboarding-d.png}\\[3pt]
    \small (г)
  \end{minipage}
  \caption{Онбординг iOS-приложения: (а)~приветственный экран с описанием приложения; (б)~выбор учебного заведения; (в)~выбор учебной группы в режиме студента; (г)~выбор преподавателя в режиме преподавателя}
  \label{fig:app-ios-onboarding}
\end{figure}

\begin{figure}[h]
  \centering
  \begin{minipage}[t]{0.45\textwidth}
    \centering
    \includegraphics[width=0.6\linewidth]{app-ios-home-a.png}\\[3pt]
    \small (а)
  \end{minipage}\hfill
  \begin{minipage}[t]{0.45\textwidth}
    \centering
    \includegraphics[width=0.6\linewidth]{app-ios-home-b.png}\\[3pt]
    \small (б)
  \end{minipage}\\[10pt]
  \begin{minipage}[t]{0.45\textwidth}
    \centering
    \includegraphics[width=0.6\linewidth]{app-ios-home-c.png}\\[3pt]
    \small (в)
  \end{minipage}\hfill
  \begin{minipage}[t]{0.45\textwidth}
    \centering
    \includegraphics[width=0.6\linewidth]{app-ios-home-d.png}\\[3pt]
    \small (г)
  \end{minipage}
  \caption{Главный экран iOS-приложения: (а)~карточка идущей пары; (б)~карточка ближайшей пары через несколько дней; (в)~контекстное меню выбора избранной учебной группы при длительном нажатии на карточку расписания; (г)~поиск справочной сущности на главном экране}
  \label{fig:app-ios-home}
\end{figure}

\begin{figure}[h]
  \centering
  \begin{minipage}[t]{0.45\textwidth}
    \centering
    \includegraphics[width=0.6\linewidth]{app-ios-entity-lists-a.png}\\[3pt]
    \small (а)
  \end{minipage}\hfill
  \begin{minipage}[t]{0.45\textwidth}
    \centering
    \includegraphics[width=0.6\linewidth]{app-ios-entity-lists-b.png}\\[3pt]
    \small (б)
  \end{minipage}\\[10pt]
  \begin{minipage}[t]{0.45\textwidth}
    \centering
    \includegraphics[width=0.6\linewidth]{app-ios-entity-lists-c.png}\\[3pt]
    \small (в)
  \end{minipage}\hfill
  \begin{minipage}[t]{0.45\textwidth}
    \centering
    \includegraphics[width=0.6\linewidth]{app-ios-entity-lists-d.png}\\[3pt]
    \small (г)
  \end{minipage}
  \caption{Списки справочных сущностей iOS-приложения: (а)~учебные группы; (б)~преподаватели; (в)~корпуса; (г)~аудитории выбранного корпуса (открываются переходом со списка корпусов)}
  \label{fig:app-ios-entity-lists}
\end{figure}

\begin{figure}[h]
  \centering
  \begin{minipage}[t]{0.45\textwidth}
    \centering
    \includegraphics[width=0.6\linewidth]{app-ios-schedule-a.png}\\[3pt]
    \small (а)
  \end{minipage}\hfill
  \begin{minipage}[t]{0.45\textwidth}
    \centering
    \includegraphics[width=0.6\linewidth]{app-ios-schedule-b.png}\\[3pt]
    \small (б)
  \end{minipage}\\[10pt]
  \begin{minipage}[t]{0.45\textwidth}
    \centering
    \includegraphics[width=0.6\linewidth]{app-ios-schedule-c.png}\\[3pt]
    \small (в)
  \end{minipage}\hfill
  \begin{minipage}[t]{0.45\textwidth}
    \centering
    \includegraphics[width=0.6\linewidth]{app-ios-schedule-d.png}\\[3pt]
    \small (г)
  \end{minipage}
  \caption{Экран расписания iOS-приложения: (а)~поденный просмотр расписания учебной группы; (б)~меню быстрой навигации по сущностям занятия при длительном нажатии на карточку; (в)~выбор произвольной даты через выдвигаемый снизу календарь; (г)~поденный просмотр расписания преподавателя}
  \label{fig:app-ios-schedule}
\end{figure}

\begin{figure}[h]
  \centering
  \begin{minipage}[t]{0.45\textwidth}
    \centering
    \includegraphics[width=0.6\linewidth]{app-ios-lesson-a.png}\\[3pt]
    \small (а)
  \end{minipage}\hfill
  \begin{minipage}[t]{0.45\textwidth}
    \centering
    \includegraphics[width=0.6\linewidth]{app-ios-lesson-b.png}\\[3pt]
    \small (б)
  \end{minipage}\\[10pt]
  \begin{minipage}[t]{0.45\textwidth}
    \centering
    \includegraphics[width=0.6\linewidth]{app-ios-lesson-c.png}\\[3pt]
    \small (в)
  \end{minipage}\hfill
  \begin{minipage}[t]{0.45\textwidth}
    \centering
    \includegraphics[width=0.6\linewidth]{app-ios-lesson-d.png}\\[3pt]
    \small (г)
  \end{minipage}
  \caption{Карточка занятия и работа с заметками iOS-приложения: (а)~карточка обычного занятия; (б)~карточка занятия с несколькими одновременно идущими дисциплинами; (в)~форма добавления заметки к занятию; (г)~экран просмотра заметки}
  \label{fig:app-ios-lesson}
\end{figure}

\begin{figure}[h]
  \centering
  \begin{minipage}[t]{0.45\textwidth}
    \centering
    \includegraphics[width=0.6\linewidth]{app-ios-settings-a.png}\\[3pt]
    \small (а)
  \end{minipage}\hfill
  \begin{minipage}[t]{0.45\textwidth}
    \centering
    \includegraphics[width=0.6\linewidth]{app-ios-settings-b.png}\\[3pt]
    \small (б)
  \end{minipage}\\[10pt]
  \begin{minipage}[t]{0.45\textwidth}
    \centering
    \includegraphics[width=0.6\linewidth]{app-ios-settings-c.png}\\[3pt]
    \small (в)
  \end{minipage}\hfill
  \begin{minipage}[t]{0.45\textwidth}
    \centering
    \includegraphics[width=0.6\linewidth]{app-ios-settings-d.png}\\[3pt]
    \small (г)
  \end{minipage}
  \caption{Настройки iOS-приложения: (а)~корневой экран настроек; (б)~выбор учебного заведения; (в)~расписание звонков (сетка пар); (г)~смена избранной учебной группы или преподавателя}
  \label{fig:app-ios-settings}
\end{figure}

\begin{figure}[h]
  \centering
  \begin{minipage}[t]{0.45\textwidth}
    \centering
    \includegraphics[width=0.6\linewidth]{app-ios-notes-favorites-cache-a.png}\\[3pt]
    \small (а)
  \end{minipage}\hfill
  \begin{minipage}[t]{0.45\textwidth}
    \centering
    \includegraphics[width=0.6\linewidth]{app-ios-notes-favorites-cache-b.png}\\[3pt]
    \small (б)
  \end{minipage}\\[10pt]
  \begin{minipage}[t]{0.45\textwidth}
    \centering
    \includegraphics[width=0.6\linewidth]{app-ios-notes-favorites-cache-c.png}\\[3pt]
    \small (в)
  \end{minipage}\hfill
  \begin{minipage}[t]{0.45\textwidth}
    \centering
    \includegraphics[width=0.6\linewidth]{app-ios-notes-favorites-cache-d.png}\\[3pt]
    \small (г)
  \end{minipage}
  \caption{Заметки, избранное и управление кэшем iOS-приложения: (а)~актуальные заметки; (б)~архив заметок; (в)~избранные преподаватели; (г)~управление локальным кэшем (раздел <<Данные и память>>)}
  \label{fig:app-ios-notes-favorites-cache}
\end{figure}

\begin{figure}[h]
  \centering
  \begin{minipage}[t]{0.45\textwidth}
    \centering
    \includegraphics[width=0.6\linewidth]{app-ios-appearance-a.png}\\[3pt]
    \small (а)
  \end{minipage}\hfill
  \begin{minipage}[t]{0.45\textwidth}
    \centering
    \includegraphics[width=0.6\linewidth]{app-ios-appearance-b.png}\\[3pt]
    \small (б)
  \end{minipage}\\[10pt]
  \begin{minipage}[t]{0.45\textwidth}
    \centering
    \includegraphics[width=0.6\linewidth]{app-ios-appearance-c.png}\\[3pt]
    \small (в)
  \end{minipage}\hfill
  \begin{minipage}[t]{0.45\textwidth}
    \centering
    \includegraphics[width=0.6\linewidth]{app-ios-appearance-d.png}\\[3pt]
    \small (г)
  \end{minipage}
  \caption{Оформление iOS-приложения: (а)~экран оформления со стандартными значениями; (б)~экран оформления с тёмной темой, иным акцентом и альтернативной иконкой; (в), (г)~примеры экранов приложения с нестандартным оформлением}
  \label{fig:app-ios-appearance}
\end{figure}

\begin{figure}[h]
  \centering
  \begin{minipage}[t]{0.45\textwidth}
    \centering
    \includegraphics[width=0.6\linewidth]{app-ios-widgets-a.png}\\[3pt]
    \small (а)
  \end{minipage}\hfill
  \begin{minipage}[t]{0.45\textwidth}
    \centering
    \includegraphics[width=0.6\linewidth]{app-ios-widgets-b.png}\\[3pt]
    \small (б)
  \end{minipage}\\[10pt]
  \begin{minipage}[t]{0.45\textwidth}
    \centering
    \includegraphics[width=0.6\linewidth]{app-ios-widgets-c.png}\\[3pt]
    \small (в)
  \end{minipage}\hfill
  \begin{minipage}[t]{0.45\textwidth}
    \centering
    \includegraphics[width=0.6\linewidth]{app-ios-widgets-d.png}\\[3pt]
    \small (г)
  \end{minipage}
  \caption{Виджеты iOS-приложения: (а)~системный лист добавления виджета с виджетом ВКурсе; (б)~рабочий стол с виджетом количества пар; (в)~рабочий стол с виджетом заметок; (г)~рабочий стол с альтернативной иконкой, тёмной темой и виджетами в согласованной оформительской палитре}
  \label{fig:app-ios-widgets}
\end{figure}
